\documentclass[aps,prd,reprint,nofootinbib,longbibliography]{revtex4-2}

\usepackage{amsmath,amssymb,amsthm,bm}
\usepackage{booktabs}
\usepackage{graphicx}
\usepackage{hyperref}
\usepackage{microtype}
\usepackage{pgfplots}
\pgfplotsset{compat=1.15,
  paperaxis/.style={
    width=\columnwidth, height=0.58\columnwidth,
    grid=major, grid style={gray!20},
    tick label style={font=\footnotesize},
    label style={font=\footnotesize},
    legend style={font=\footnotesize, draw=none, fill=none},
    line width=0.9pt}}

\newtheorem{definition}{Definition}
\newcommand{\Tr}{\operatorname{Tr}}
\newcommand{\Ac}{\mathcal{A}}
\newcommand{\Nc}{\mathcal{N}}
\newcommand{\Dc}{\mathcal{D}}
\newcommand{\Ec}{\mathcal{E}}

\begin{document}

\title{Consumer-Relative Quantum Distinguishability on Causal Boundaries}

\author{Andrew H. Bond}
\email{andrew.bond@sjsu.edu}
\thanks{ORCID 0009-0003-2599-6158}
\affiliation{San Jos\'e State University, San Jos\'e, California, USA}
\date{August 4, 2026}

\begin{abstract}
Distinguishability between quantum states is usually treated as a property of
the states alone.  We treat it instead as a property of the states together
with a declared consumer, meaning the agent whose access channel, task, and
budget determine which distinctions can be claimed.  We define the consumer
channel, the consumer-relative divergence, consumer equivalence, and
task-relative distortion, verify that the resulting instruments reproduce
exact classical anchors from an earlier projection-entropy campaign, and
measure their behavior in a small spin-chain model where every quantity is
computable.  Three findings result.  First, a nested hierarchy of consumers
divides one fixed distinguishability total in a way that moves with the
excitation, so a certificate built from one observable class can be exactly
vacuous for an excitation in another class.  Second, at scarce storage
budgets an encoder optimized for a declared energy functional preserves that
functional exactly while losing badly on global state fidelity, the encoder
optimized for fidelity does the opposite, and the trade disappears once the
budget is large enough to store everything the task needs.  Third, a
consistency requirement imposed across a family of consumers does no work
when the consumers are probed disjointly and does real work, selecting the
weak-coupling regime of the model family, exactly when all consumers witness
one shared event.  We also report three instrument defects found during
construction and converted into standing controls, because each is a silent
failure mode of restricted-divergence numerics.  Everything numerical in
this paper is demonstrated in finite-dimensional models.  Nothing in it is a
claim about physical gravity, real horizons, or quantum gravity, and the
consumer does not create anything physical by choosing what to observe.
\end{abstract}

\maketitle

\section{Introduction}
\label{sec:intro}

Several modern derivations of gravitational field equations from information
pass through a single kind of object, the distinguishability of quantum
states as seen from a restricted vantage point.  Jacobson derived the
Einstein equation as an equation of state by demanding that a thermodynamic
identity hold for every local causal horizon, which is a universal
quantifier over restricted observers \cite{Jacobson1995}.  Bianconi
postulated an action that is itself a quantum relative entropy between
geometric and matter degrees of freedom \cite{Bianconi2025}.  Dorau and Much
proved, for a special class of states on a special geometry, that the
relative entropy accessible to an observer confined outside a horizon equals
an energy flux, and recovered semiclassical gravity from that equality
together with an assumed entropy-area formula \cite{DorauMuch2026}.  A
related construction runs the thermodynamic argument on stretched light
cones \cite{AlonsoSerrano2025}.  In each case the restriction to an
observer's accessible algebra is the informational move, and in each case
the observer appears only as a restriction.  None of these constructions
contains a consumer in the operational sense, an agent with a declared task
whose success is scored on that task rather than on full state
reconstruction.

This paper supplies and tests that missing layer in the only setting where
every quantity can be computed and audited, namely finite-dimensional
models.  The program it extends, called Observation Theory in our earlier
classical work, holds that distinguishability is not one thing.  It is
relative to a consumer, and the distinctions a consumer needs for its task
can be preserved while global state fidelity is sacrificed, and conversely.
The classical version of that statement was established with projection and
entropy instruments in prior work of this campaign.  Here the classical read
operator is replaced by a quantum channel, the most general physically
allowed map on quantum states, and the same questions are asked again.

Three claim types are kept separate throughout.  Instrument claims state
that the implemented objects compute what the definitions name, and they are
verified against closed forms.  In-model structural claims state measured
facts about one declared family of finite-dimensional models.  Physical
claims about gravity, real horizons, or quantum gravity are absent, and no
sentence of this paper should be read as making one.  The causal direction
in every construction runs from the physical state and the model couplings,
through the accessible channel, to the consumer-relative divergence.  No
choice of channel, task, or budget feeds back into the state or the
couplings.  The consumer does not create anything physical by choosing what
to observe.  Observation relativity, as used here, concerns evidential
access, not reality.

One further discipline governs the whole paper.  Renaming known quantities
adds nothing.  Calling relative entropy distinguishability, calling the
data-processing inequality an access law, and calling a partial trace a
consumer channel are relabelings of standard mathematics.  The contribution,
if any, is confined to the task layer defined below and to the measured
results.  Where a result turns out to be a relabeling, the paper says so.

The paper proceeds in six steps.  It replaces the classical read operator
with a quantum channel and states the access law.  It defines consumer
equivalence and task-relative distortion.  It verifies exact classical
reduction against anchors inherited from the earlier campaign.  It analyzes
a small transverse-field Ising model with a declared inside and outside,
reporting instrument results and three defects converted to standing
controls.  It measures a consumer hierarchy, a budgeted encoding trade we
call the flip, and a consumer-family consistency requirement with a
two-sided outcome.  It closes by stating precisely what would be required to
connect any of this to the gravitational results cited above, and why the
present constructions do not.

\section{Definitions}
\label{sec:defs}

\subsection{Consumer channel and consumer-relative divergence}

A quantum channel is a completely positive trace-preserving linear map
$\Nc$ on density matrices \cite{NielsenChuang}.  A consumer $C$ receives
states only through its declared channel $\Nc_C$.  Canonical instances used
in this paper are the partial trace over an inaccessible tensor factor, local
dephasing at declared strength, and pinching, meaning full dephasing in the
computational basis, which restricts the consumer to the diagonal algebra.

The default divergence is the Umegaki quantum relative entropy
\cite{Umegaki1962},
\begin{equation}
D(\rho\,\|\,\sigma) \;=\; \Tr\!\left[\rho\left(\ln\rho-\ln\sigma\right)\right],
\label{eq:umegaki}
\end{equation}
in nats, the natural-logarithm unit of information, defined as $+\infty$
when the support of $\rho$ is not contained in the support of $\sigma$.  The consumer-relative divergence is
\begin{equation}
D_C(\rho\,\|\,\sigma) \;=\; D\!\left(\Nc_C(\rho)\,\|\,\Nc_C(\sigma)\right).
\label{eq:dc}
\end{equation}

\subsection{The access law}

For every channel $\Nc$ and all states $\rho,\sigma$,
\begin{equation}
D\!\left(\Nc(\rho)\,\|\,\Nc(\sigma)\right) \;\le\; D(\rho\,\|\,\sigma),
\label{eq:dpi}
\end{equation}
the data-processing inequality \cite{Lindblad1975}.  Equality
characterizes sufficiency of the channel for the pair \cite{Petz1986}.  In
consumer language Eq.~\eqref{eq:dpi} says that claimable distinctions cannot
exceed underlying distinctions, and that restriction, coarse-graining, and
inaccessibility only erase distinguishability, never create it.  This
reading is a relabeling of a known theorem.  The instruments of
Sec.~\ref{sec:model} do not test whether Eq.~\eqref{eq:dpi} is true.  They
verify that the implementation computes the objects the theorem is about,
which Sec.~\ref{sec:model} shows is a real failure surface.

\subsection{Consumer equivalence}

Let $\Ac_C$ be the algebra of observables the consumer can evaluate on the
output of its channel.  Two states are consumer equivalent,
\begin{equation}
\rho \sim_C \sigma
\quad\Longleftrightarrow\quad
\Tr(\rho\,O)=\Tr(\sigma\,O)\ \ \text{for all } O\in\Ac_C ,
\label{eq:equiv}
\end{equation}
when no observable available to the consumer separates them.  The quotient
of state space by $\sim_C$ is the world the consumer can act on, with
unreadable distinctions collapsed.  In this paper the channel is the
primitive object and $\Ac_C$ is derived from it, because two consumers with
the same algebra can differ in noise, and a conditional expectation onto a
subalgebra exists only for special algebra-state pairs.

\subsection{Task-relative distortion}

A task is a declared functional $f$ on states, fixed before any sweep.  For
an encoding map $\Ec$ at a declared budget, with a declared decoder $\Dc$,
the task distortion is the error in the functional,
\begin{equation}
d_f(\rho,\Ec) \;=\; \bigl|\,f(\rho)-f\!\left(\Dc(\Ec(\rho))\right)\bigr| ,
\label{eq:taskdist}
\end{equation}
scored against the global reconstruction error
\begin{equation}
1-F\!\left(\rho,\Dc(\Ec(\rho))\right),
\qquad
F(\rho,\sigma)=\Bigl(\Tr\sqrt{\sqrt{\rho}\,\sigma\sqrt{\rho}}\Bigr)^{2},
\label{eq:infid}
\end{equation}
where $F$ is the Uhlmann fidelity \cite{Uhlmann1976,Jozsa1994}.  The task
layer is where the paper's contribution lives, if anywhere.  The divergence
layer of Eqs.~\eqref{eq:umegaki} through \eqref{eq:equiv} answers how
distinguishable two states are through a channel.  The task layer asks which
distinction a particular consumer needs, and the experiments below measure
whether that question has content beyond the divergence layer.

\subsection{Boundary language}

Throughout, a boundary is a statement about which tensor factor of a
finite-dimensional Hilbert space a consumer can touch, with the consumer
channel a partial trace over the rest.  A bipartite lattice cut is already a
boundary in this sense.  The causal boundaries of the gravitational
literature are the special case in which the split is causal and the
restricted state is thermal with respect to a modular flow, the intrinsic
evolution that the state itself induces on the accessible algebra.  The models of
this paper use the general case and claim nothing about the special one.
Section~\ref{sec:gap} lists what the special case would additionally
require.

\section{Classical reduction and anchors}
\label{sec:classical}

When all states and channels are simultaneously diagonalizable, the quantum
objects must collapse to the classical ones.  Equation~\eqref{eq:umegaki}
becomes the Kullback-Leibler divergence, the channel becomes a stochastic
map, and consumer equivalence becomes the classical quotient.  The earlier
campaign supplies two exact anchors, the branch-weight distributions of two
fold projections whose entropies are known in closed form.  A two-branch
fold carries branch weights $(\tfrac12,\tfrac12)$ and exactly one bit.  A
double fold carries branch weights $(\tfrac14,\tfrac12,\tfrac14)$ at its
symmetric slice and exactly $1.5$ bits.

The declared embedding, fixed before any run, maps fold branches in
increasing order of the evolution parameter to computational basis states,
coarea weights to eigenvalues of a diagonal density matrix, the uniform
distribution to the reference state $I/d$, and the branch-label algebra to
the diagonal algebra.  Entropy is then recovered through the divergence
pathway,
\begin{equation}
S(\rho) \;=\; \log_2 d \;-\; \frac{D(\rho\,\|\,I/d)}{\ln 2}.
\label{eq:entropy}
\end{equation}
The instrument reproduces $1.000000000000$ bits for the two-branch fold and
$1.500000000000$ bits for the double fold through
Eq.~\eqref{eq:entropy}, at residuals below $2\times 10^{-15}$.

The anchors also pin the equality case of the access law.  Merging the two
orientation branches of the double fold is a classical channel.  Against
the uniform reference the merge loses exactly nothing, because the
conditional distributions inside the merged group coincide, so the
data-processing gap is an exact sufficiency equality, measured at
$1.1\times 10^{-16}$.  Against the asymmetric reference
$(\tfrac12,\tfrac14,\tfrac14)$ the same merge loses a nonzero amount with
the closed form
\begin{equation}
\Delta \;=\; \frac{\ln 2}{4}-\frac{1}{2}\ln\frac{4}{3}
\;\approx\; 0.029446\ \text{nats},
\label{eq:gap}
\end{equation}
and the measured gap matches Eq.~\eqref{eq:gap} to twelve decimals.  A
quantum instrument that failed these classical limits would be invalid
regardless of anything else it computed.  The embedding is declared rather
than unique, and the reduction check is exactly as strong as that
declaration.

\section{Model and instrument}
\label{sec:model}

\subsection{Declared model}

The model for all quantum runs is an open chain of $8$ qubits with the
transverse-field Ising Hamiltonian
\begin{equation}
H \;=\; -J\sum_{k=0}^{6} Z_k Z_{k+1} \;-\; h\sum_{k=0}^{7} X_k ,
\label{eq:ising}
\end{equation}
in the gapped paramagnetic regime.  The reference state is thermal,
$\sigma_\beta = e^{-\beta H}/\Tr\,e^{-\beta H}$, at $\beta=0.5$ with $J=1$
and $h=2$ for the instrument, hierarchy, and flip studies, and at
$\beta=0.4$ over the coupling grid $J\in\{0.6,0.8,1.0,1.2\}$,
$h\in\{1.6,1.8,2.0,2.2\}$ for the family study of Sec.~\ref{sec:family}.
Excitations are local unitaries $U_\theta=\exp(i\theta P_s)$ for a
single-site Pauli operator $P_s$, applied to the thermal reference.  The
declared cut places sites $0$ through $3$ inside and sites $4$ through $7$
outside, and the exterior consumer's channel is the partial trace onto the
outside.  All computations are exact diagonalization in double precision.

\subsection{Instrument results}

The instrument layer verifies that the implementation computes
Eqs.~\eqref{eq:umegaki} through \eqref{eq:dc} correctly.  Closed-form
checks cover the commuting Kullback-Leibler case, a pure state against the
maximally mixed state, self-divergence, unitary invariance, and the
support-violation convention, with a maximum residual below
$2\times 10^{-15}$.  Across the excitation and channel sweep the
implementation performed $32$ data-processing checks of
Eq.~\eqref{eq:dpi}, every margin strictly positive and finite, minimum
margin $9.6\times 10^{-8}$.  Composition monotonicity, meaning that
composing two channels never increases the divergence, held at every point.
Sweeping the numerical support floor across
$10^{-14}$ through $10^{-10}$ changed no reported divergence, with a
measured spread of zero.

\subsection{Three defects converted to standing controls}
\label{sec:defects}

Three instrument defects were discovered by failed first designs.  Each is
reported here as a methods result, because each is a silent failure mode of
restricted-divergence numerics that produces passing-looking output, and
each is now a standing control asserted on every run.

\emph{Silent infinite divergence at the support floor.}
At $\beta=1$ the smallest thermal eigenvalue, near $2\times 10^{-12}$, sat
at the $10^{-12}$ support floor of the divergence routine.  A full-rank
reference was thereby misread as rank deficient, the full-state divergence
went silently infinite, and every global-versus-restricted margin became
infinity minus a finite number.  All margins were vacuously positive, and
the defect was invisible in the minimum margin.  The production model must
now clear the support floor by three decades, which is asserted, and every
global divergence carries a finiteness assertion.  The first evidence
record contained this flaw and was regenerated.  The corrected record
supersedes it.

\emph{No-signalling null.}
A unitary, or any channel, acting only on the traced-out interior cannot
change the exterior reduced state at all.  The first sweep design placed
the excitation inside, so every consumer divergence was exactly zero and
the data-processing checks passed trivially on a degenerate sweep.
Interior excitation is now an asserted null control, with the measured
exterior divergence below $7\times 10^{-16}$ for excitations both deep
inside and at the interior side of the cut.

\emph{Symmetry null.}
By the global spin-flip symmetry of Eq.~\eqref{eq:ising}, the single-site
reduced thermal state is exactly $(I+mX)/2$ for a magnetization parameter
$m$.  An $X$-axis rotation commutes with it and is invisible to the
site-only consumer even though it acts on that consumer's own qubit, with
the measured on-site divergence below $7\times 10^{-16}$.  Whether a
consumer can see an excitation depends on the interplay of excitation
direction and state symmetry, not only on locality.  This null is also the
first concrete instance in this paper of certificate vacuity, developed
quantitatively in Sec.~\ref{sec:hierarchy}.

\begin{figure}[t]
\centering
\begin{tikzpicture}
\begin{axis}[paperaxis, height=0.52\columnwidth,
  ybar=2pt, bar width=7pt,
  xtick={1,2,3,4},
  xticklabels={position, position$+$flux, boundary, full state},
  x tick label style={font=\footnotesize, rotate=20, anchor=north east},
  ylabel={$D_{C}$ (nats)}, ymin=0,
  legend pos=north west, enlarge x limits=0.12]
\addplot[fill=blue!65!black!60, draw=blue!70!black]
  table[x=index, y=zchain] {figdata/qo1_hierarchy.dat};
\addlegendentry{$Z$ excitation}
\addplot[fill=red!65!black!45, draw=red!70!black]
  table[x=index, y=xchain] {figdata/qo1_hierarchy.dat};
\addlegendentry{$X$ excitation}
\end{axis}
\end{tikzpicture}
\caption{Consumer hierarchy at excitation strength $0.8$.  For the $Z$
excitation the position consumer sees exactly nothing and almost the whole
distinguishability arrives at the flux rung.  For the $X$ excitation the
position consumer already sees the event through correlations and the
largest gap sits at the interior rung.  The same four consumers divide the
two distinguishability totals entirely differently.}
\label{fig:hierarchy}
\end{figure}

\section{Consumer hierarchy}
\label{sec:hierarchy}

\subsection{Design}

Four consumers are arranged as a degradation chain, each weaker consumer a
channel composed onto the stronger one's output, so the ordering of their
divergences is a theorem and any measured inversion would be an
implementation bug.  From weakest to strongest the levels are position
only, realized as a $Z$-basis pinch on every outside site, position plus
flux, which keeps the full single-site algebra of the excited site and
pinches the rest, the full boundary, which is the partial trace onto the
outside region, and the full state.  The declared flux stand-in is the
transverse-field energy density on the excited site,
\begin{equation}
f(\rho) \;=\; \Tr\!\left[\rho\,(-h\,X_{s})\right],
\label{eq:flux}
\end{equation}
declared rather than derived, a choice whose risks are discussed with the
results.  Two excitation axes are measured, a $Z$-axis rotation and an
$X$-axis rotation on the same outside site, both at
$\theta\in\{0.4,0.8,1.2\}$.

\begin{table}[t]
\caption{Consumer-hierarchy divergence chains at excitation strength
$\theta=0.8$, in nats.  Each level is a channel degradation of the one
below it, so each column is nondecreasing downward by
Eq.~\eqref{eq:dpi}.  The orderings and gap locations are stable across
$\theta=0.4$ and $\theta=1.2$.}
\label{tab:hierarchy}
\begin{ruledtabular}
\begin{tabular}{lcc}
Consumer level & $Z$ excitation & $X$ excitation \\
\colrule
position only        & $0.000000$ & $0.047621$ \\
position plus flux   & $0.667985$ & $0.116320$ \\
full boundary        & $0.675183$ & $0.148405$ \\
full state           & $0.723532$ & $0.300990$ \\
\end{tabular}
\end{ruledtabular}
\end{table}

\subsection{Results}

Table~\ref{tab:hierarchy} shows the measured chains.  For the $Z$-axis
excitation the position-only consumer sees exactly nothing.  This is a
commutation theorem, asserted in the instrument, because the excitation
commutes with every $Z$-string the pinched consumer can read.  The largest
gap, $0.667985$ nats, sits at the flux rung.  For the $X$-axis excitation
the position-only consumer does see the event, at $0.047621$ nats, because
the excitation disturbs position correlations even though the on-site
marginal is unchanged by the symmetry null of Sec.~\ref{sec:defects}.  The
largest gap for this excitation, $0.152585$ nats, sits at the interior
rung between the full boundary and the full state.

The structural statement is that the same four consumers, ordered the same
way by Eq.~\eqref{eq:dpi}, divide the same kind of distinguishability total
entirely differently depending on the excitation sector.  Where the
distinguishing information lives moves with the excitation.  A certificate
of the form ``the position consumer sees nothing, therefore nothing
happened'' is exactly vacuous for a flux-sector excitation, since the
position consumer's divergence is identically zero there while the flux
consumer's divergence is large.  This is the task layer measured rather
than relabeled.  The divergence layer alone fixes only the ordering.  The
gaps, and where they sit, are model facts that no renaming produces.  The
caveat is adjacent.  The flux functional Eq.~\eqref{eq:flux} is a declared
stand-in on a lattice with no null direction, and every statement above is
a statement about this model class at these parameters.

\begin{figure}[t]
\centering
\begin{tikzpicture}
\begin{axis}[paperaxis, height=0.55\columnwidth,
  xlabel={budget (stored qubits)}, xtick={1,2,3},
  ylabel={held-out distortion}, ymin=-0.04, ymax=1.42,
  legend style={font=\footnotesize, draw=none, fill=none,
    at={(0.5,1.02)}, anchor=south, legend columns=2}]
\begin{pgfonlayer}{axis background}
\fill[blue!6] (axis cs:0.75,-0.04) rectangle (axis cs:2.25,1.42);
\end{pgfonlayer}
\node[font=\footnotesize, gray] at (axis cs:1.5,1.05) {flip region};
\addplot[blue!70!black, mark=*]
  table[x=budget, y=Ttask] {figdata/qo2_flip.dat};
\addlegendentry{task arm, task error}
\addplot[red!70!black, mark=square*]
  table[x=budget, y=Ftask] {figdata/qo2_flip.dat};
\addlegendentry{fidelity arm, task error}
\addplot[blue!70!black, dashed, mark=o, mark options={solid}]
  table[x=budget, y=Tinfid] {figdata/qo2_flip.dat};
\addlegendentry{task arm, infidelity}
\addplot[red!70!black, dashed, mark=square, mark options={solid}]
  table[x=budget, y=Finfid] {figdata/qo2_flip.dat};
\addlegendentry{fidelity arm, infidelity}
\addplot[black!50, dotted, mark=triangle*, mark options={solid}]
  table[x=budget, y=antitask] {figdata/qo2_flip.dat};
\addlegendentry{anti-arm, task error}
\end{axis}
\end{tikzpicture}
\caption{The flip on held-out states.  At budgets one and two the task arm
preserves the flux functional exactly while paying a large infidelity
price, and the fidelity arm does the opposite.  At budget three the
fidelity-optimal subset happens to contain the flux qubit and the trade
disappears, so the flip is a scarcity phenomenon.  The anti-arm is worst
on the task at every budget, which certifies that the declared functional
governs the task.}
\label{fig:flip}
\end{figure}

\section{The flip}
\label{sec:flip}

\subsection{Design}

The flip hypothesis, transplanted from the classical program, is that at a
fixed budget an encoder optimized to preserve a declared task functional
beats a fidelity-optimized encoder on the task while losing on global
fidelity.  The declared design is as follows.  The model is the outside
region of Sec.~\ref{sec:model}, four qubits with the thermal marginal as
vacuum.  Scenario states apply a strong $Z$ rotation on a
fidelity-dominant site and a weak $Z$ rotation on the flux site, so the
fidelity signal and the task signal live on different qubits.  The task is
the flux functional Eq.~\eqref{eq:flux} on the flux site.  The encoder
family, identical for every arm, keeps a subset of $k$ of the $4$ qubits
and decodes by reattaching the vacuum marginals of the discarded qubits.
The budget is the number of stored qubit slots.  The family is discrete
and exhaustively searched, so there is no optimizer whose artifacts could
manufacture the result.  Arms are selected on a training grid of scenario
angles and judged only on a disjoint held-out grid, with declared win
margins of $0.01$ on both axes.  The F-arm minimizes mean infidelity, the
T-arm minimizes mean task distortion, and a mandatory anti-arm applies the
T-arm's encoder with the flux qubit dephased before storage, spending the
identical budget while destroying the flux observable.  If the anti-arm
were not clearly worst on the task, the declared functional would not be
what governs the task and the run would be uninterpretable rather than
positive.  A two-qubit analytic flip control, which no search could fake,
is part of the accompanying test suite.

\begin{table}[t]
\caption{Held-out flip results by budget $k$ in stored qubit slots.  Task
distortion is Eq.~\eqref{eq:taskdist} for the flux functional, infidelity
is Eq.~\eqref{eq:infid}.  Task values of $0.000$ are exact zeros of the
encoding in exact arithmetic, measured at $10^{-16}$.}
\label{tab:flip}
\begin{ruledtabular}
\begin{tabular}{clccc}
$k$ & Arm & Keeps & Task distortion & Infidelity \\
\colrule
1 & F    & strong site        & $0.118$ & $0.069$ \\
1 & T    & flux site          & $0.000$ & $0.326$ \\
1 & anti & flux site, pinched & $1.289$ & $0.431$ \\
\colrule
2 & F    & strong pair        & $0.118$ & $0.046$ \\
2 & T    & flux pair          & $0.000$ & $0.330$ \\
2 & anti & flux pair, pinched & $1.289$ & $0.430$ \\
\colrule
3 & F    & includes flux site & $0.000$ & $0.025$ \\
3 & T    & includes flux site & $0.000$ & $0.050$ \\
3 & anti & flux site pinched  & $1.289$ & $0.197$ \\
\end{tabular}
\end{ruledtabular}
\end{table}

\subsection{Results}

Table~\ref{tab:flip} reports the held-out outcomes.  At budget $k=1$ the
T-arm keeps the flux site and achieves task distortion exactly zero,
because a kept single-site marginal determines the single-site functional
exactly, at the price of infidelity $0.326$.  The F-arm keeps the strong
site and achieves infidelity $0.069$ at task distortion $0.118$.  The
task-optimized encoder wins the task outright and is worse on global
fidelity by a factor of $4.7$.  This is the flip.  At $k=2$ the split
persists, and the fidelity arm's choice is informative in its own right,
since it prefers the correlated strong pair over adding the weakly excited
flux qubit.  Fidelity favors preserving correlation structure.  At $k=3$
the fidelity-optimal subset happens to contain the flux qubit, both arms
reach task distortion zero, and the flip vanishes.  The anti-arm is worst
on the task at every budget, $1.289$ against at most $0.118$, so the
declared functional does govern the task and the runs are interpretable.

The flip region is contiguous, budgets $1$ and $2$, and the flip dies at
exactly the budget where storage suffices to keep everything
task-relevant.  The flip is a scarcity phenomenon.  This is the
budget-relativity structure the classical program predicts, here realized
with quantum states, a quantum channel budget, and an energy functional.
The result is not an estimation protocol.  Classical shadow methods
predict many observables of a state from few randomized measurements with
guarantees that sharpen as the observable class narrows
\cite{Huang2020}.  The flip instead measures a strict trade between one
declared functional and global reconstruction at matched storage, and its
content is the existence and location of that trade, not an estimator.
Two caveats are adjacent.  The scenario family was declared so that task
signal and fidelity signal live on different qubits, which is the
misalignment the hypothesis requires, so the result demonstrates
existence in a declared model class and nothing stronger.  These runs
carry an exploratory label, and a sealed preregistered replication would
be required before the flip could be claimed beyond this model class.

\begin{figure}[t]
\centering
\begin{tikzpicture}
\begin{axis}[paperaxis, height=0.5\columnwidth,
  xlabel={family size}, xtick={1,2,3,4,5,6},
  ylabel={joint survival fraction}, ymin=0, ymax=1.08,
  legend pos=south west, unbounded coords=discard]
\addplot[black!60, mark=square*]
  table[x=familysize, y=jointA] {figdata/qo3_survival.dat};
\addlegendentry{family A, disjoint probes}
\addplot[blue!70!black, mark=*]
  table[x=familysize, y=jointB] {figdata/qo3_survival.dat};
\addlegendentry{family B, shared event}
\end{axis}
\end{tikzpicture}
\caption{Fraction of the coupling grid on which the family requirement
holds, as consumers are added.  For disjoint probes the joint requirement
never removes a coupling point that a single consumer keeps, so the curve
stays at one and the quantifier is inert.  For nested windows witnessing
one shared event the joint requirement tightens to one quarter of the
grid, strictly below every single-consumer value, and the surviving
quarter is exactly the weak-coupling row.}
\label{fig:family}
\end{figure}

\section{Consumer-family consistency}
\label{sec:family}

\subsection{Design}

Jacobson's derivation runs on a universal quantifier, a demand imposed for
every local horizon at once \cite{Jacobson1995}.  The model-internal
question tested here is whether such a quantifier does any work, meaning
whether a requirement imposed across a whole family of consumers
constrains a model more than the same requirement imposed on each consumer
alone.  The declared requirement, per coupling point $(J,h)$ of the model
family at $\beta=0.4$, is that each consumer's divergence follow the
declared flux weight,
\begin{equation}
D_c(\theta) \;=\; \lambda_c\,\theta^{2}
\quad\text{within relative residual } \epsilon,
\label{eq:law}
\end{equation}
with the constant $\lambda_c$ universal across the family within spread
$\delta$, where the spread is the range of the constants divided by their
mean.  The flux weight $\theta^{2}$ is declared from the excitation
parameter alone, so nothing informational enters the declaration and the
universality content is measured rather than inserted.  Thresholds
$\epsilon=\delta=0.05$ were fixed with the design, and all verdicts are
also reported at a loose value $0.15$.  Excitation strengths are
$\theta\in\{0.05,0.10,0.15,0.20\}$.

Two families are measured.  Family A consists of six wedge consumers, the
right exteriors of cuts $1$ through $6$, each probed by a $Z$ rotation at
its own boundary site, so the probes are disjoint.  Family B consists of
five nested windows, from the single shared site up to the full chain, all
witnessing one shared excitation at one site, so every consumer sees the
same event with a different extent of access.

\subsection{Family A, disjoint probes}

The quantifier is inert.  The $\theta^{2}$ law holds for every consumer at
every coupling point, the constants $\lambda_c$ agree across the six
consumers within a relative spread of at most $1.6\times 10^{-3}$ over the
whole grid, more than an order of magnitude inside the declared $\delta$,
and the joint requirement never removes a coupling point that any single
consumer keeps.  Survival is $1.0$ for
every consumer and for every subfamily at both thresholds.  The mechanism
is visible in the construction.  At short correlation length the effect of
a boundary-local probe is confined near its site up to tails set by that
length, every wedge containing its own probe captures it essentially
completely, and disjoint probes tie nothing together.  A
family of observers probed disjointly cannot make an all-observers
quantifier do work.

\begin{table}[t]
\caption{Family B universality spread of the constants $\lambda$ across
the five nested windows, by coupling point.  The strict threshold is
$\delta=0.05$.  The spread increases monotonically with the coupling
ratio $J/h$, which sets the correlation strength.}
\label{tab:spread}
\begin{ruledtabular}
\begin{tabular}{ccccc}
 & $h=1.6$ & $h=1.8$ & $h=2.0$ & $h=2.2$ \\
\colrule
$J=0.6$ & $0.0350$ & $0.0345$ & $0.0340$ & $0.0335$ \\
$J=0.8$ & $0.0612$ & $0.0603$ & $0.0594$ & $0.0585$ \\
$J=1.0$ & $0.0936$ & $0.0922$ & $0.0908$ & $0.0894$ \\
$J=1.2$ & $0.1313$ & $0.1293$ & $0.1273$ & $0.1253$ \\
\end{tabular}
\end{ruledtabular}
\end{table}

\subsection{Family B, one shared event}

The quantifier is active.  Every window passes the $\theta^{2}$ law singly
at every coupling point, so no single consumer constrains the coupling
grid at all.  The joint universality requirement tightens strictly with
family size, with surviving fractions of the sixteen-point coupling grid
of $1.00$, $0.75$, and $0.25$ at family sizes one, two, and three, and
$0.25$ thereafter.  The joint requirement is strictly tighter than every
single-consumer requirement.  The surviving quarter is exactly the
weak-coupling row $J=0.6$.  Table~\ref{tab:spread} shows why.  The
universality spread orders monotonically with correlation strength, from
$0.0335$ at the weakest-coupling corner to $0.1313$ at the
strongest-coupling corner, so the family requirement constrains couplings
by bounding correlation assistance and selects the near-factorized regime.
The constants also saturate once the window clears the correlation
length.  On the surviving row the three-site window and every larger
window agree to four decimals, while at the strongest coupling that
agreement is reached only by the two largest windows.  All sufficiently
large observers agree, and only observers smaller than the correlation
length dissent.

\subsection{Reading}

Stated carefully and model-internally, an all-observers consistency
requirement does real constraining work precisely when the observers
overlap on shared events, and does none when they are probed disjointly.
The overlap structure is what the continuum family of local horizons in
Jacobson's argument possesses and what family A lacks.  This is a
statement about one finite model family and about the structure of the
quantifier, not about gravity.  Two caveats are adjacent.  The verdict is
threshold-relative, since at the loose threshold of $0.15$ every coupling
point survives every requirement, so a sealed study must preregister the
spread threshold before it can claim a selection effect.  And in one
dimension every cut has a boundary of a single site, so this experiment
tests the universality content of the quantifier and says nothing about
area scaling.

\section{What a connection to gravity would require}
\label{sec:gap}

This section states precisely what separates the constructions above from
the gravitational results of
Refs.~\cite{Jacobson1995,DorauMuch2026}, so that no reader mistakes the
finite models for more than they are.

Every algebra in this paper is a finite-dimensional type I factor.  The
observable algebra of a real horizon exterior is a type III von Neumann
factor, which admits no density matrices, no partial trace, and no von
Neumann entropy of the restricted state.  Connecting the consumer layer to
that setting would require, at minimum, four replacements.  First, type
III algebras in place of finite factors, with relative entropy defined
through relative modular operators rather than
Eq.~\eqref{eq:umegaki}.  Second, a genuine causal structure in place of a
lattice cut, since the cut used here has no light cones, no Killing flow,
and no notion of a horizon generator.  Third, the entropy-area input
justified rather than declared, because in Refs.~\cite{Jacobson1995,
DorauMuch2026} the proportionality of entropy to area is an assumption,
and nothing in this paper produces or replaces it.  Fourth, the special
state class relaxed, since the theorem of Ref.~\cite{DorauMuch2026} holds
for coherent excitations on a bifurcate Killing horizon and the thermal
lattice states used here have no counterpart standing.

The finite-dimensional statements divide accordingly.  The existence of
reduced density matrices, of the pinching channel, and of
Eq.~\eqref{eq:entropy} are artifacts of type I, and every number in
Secs.~\ref{sec:model} through \ref{sec:family} depends on them.  The
relative entropy itself and the data-processing inequality have exact
type III counterparts, so the definitions of Sec.~\ref{sec:defs} survive
the passage in principle.  Whether the task layer survives it, meaning
whether consumer equivalence for declared functionals and the flip can
even be formulated without density matrices, is open, and this paper
takes no position on it.

\section{Conclusion}
\label{sec:conclusion}

This paper replaced the classical read operator of an earlier
observation-theoretic program with a quantum channel, defined
consumer-relative divergence, consumer equivalence, and task-relative
distortion, and measured all of them in finite-dimensional models with
exact classical anchors.  The instrument layer reproduced closed forms at
the $10^{-15}$ level, held every data-processing margin strictly
positive, and yielded three transferable methods findings, the
support-floor infinity, the no-signalling null, and the symmetry null,
each a silent failure mode of restricted-divergence numerics now held as
a standing control.  The hierarchy result showed that where
distinguishability lives moves with the excitation sector, making some
certificates exactly vacuous.  The flip result showed a strict
task-versus-fidelity trade that exists at scarce budgets and dies when
scarcity ends.  The family result showed that an all-observers
requirement constrains a model exactly when the observers share events,
and is otherwise inert.

The relabeling clause applies to everything else.  Whatever in this paper
merely renames relative entropy, data processing, or the partial trace is
not a contribution.  The contribution is confined to the task layer and
to the measured results, all of which are demonstrated in
finite-dimensional models, carry an exploratory evidence label pending
sealed preregistered replication, and make no claim about physical
gravity, real horizons, or quantum gravity.  The consumer does not create
anything physical by choosing what to observe.

\begin{acknowledgments}
The author thanks the maintainers of the open-source scientific Python
stack used for all computations.
\end{acknowledgments}

\begin{thebibliography}{11}

\bibitem{Umegaki1962}
H.~Umegaki,
``Conditional expectation in an operator algebra, IV (entropy and
information),''
K\=odai Math. Sem. Rep. \textbf{14}, 59 (1962).

\bibitem{Lindblad1975}
G.~Lindblad,
``Completely positive maps and entropy inequalities,''
Commun. Math. Phys. \textbf{40}, 147 (1975).

\bibitem{Uhlmann1976}
A.~Uhlmann,
``The `transition probability' in the state space of a $^{*}$-algebra,''
Rep. Math. Phys. \textbf{9}, 273 (1976).

\bibitem{Jozsa1994}
R.~Jozsa,
``Fidelity for mixed quantum states,''
J. Mod. Opt. \textbf{41}, 2315 (1994).

\bibitem{Petz1986}
D.~Petz,
``Sufficient subalgebras and the relative entropy of states of a von
Neumann algebra,''
Commun. Math. Phys. \textbf{105}, 123 (1986).

\bibitem{NielsenChuang}
M.~A.~Nielsen and I.~L.~Chuang,
\emph{Quantum Computation and Quantum Information}
(Cambridge University Press, Cambridge, 2010).

\bibitem{Jacobson1995}
T.~Jacobson,
``Thermodynamics of spacetime: The Einstein equation of state,''
Phys. Rev. Lett. \textbf{75}, 1260 (1995),
arXiv:gr-qc/9504004.

\bibitem{Bianconi2025}
G.~Bianconi,
``Gravity from entropy,''
Phys. Rev. D \textbf{111}, 066001 (2025),
arXiv:2408.14391.

\bibitem{DorauMuch2026}
P.~Dorau and A.~Much,
``From quantum relative entropy to the semiclassical Einstein
equations,''
Phys. Rev. Lett. \textbf{136}, 091602 (2026),
arXiv:2510.24491.

\bibitem{AlonsoSerrano2025}
A.~Alonso-Serrano, L.~J.~Garay, M.~Liska, and C.~Lopez~Pineros,
``Gravity from equilibrium thermodynamics of stretched light cones,''
Phys. Rev. D \textbf{112}, 124026 (2025),
arXiv:2509.08566.

\bibitem{Huang2020}
H.-Y.~Huang, R.~Kueng, and J.~Preskill,
``Predicting many properties of a quantum system from very few
measurements,''
Nat. Phys. \textbf{16}, 1050 (2020).

\end{thebibliography}

\end{document}
